% =============================================================
%  PROGRESS REPORT  (10-page body)
%  Kaustubh Pandey (2025PA17326) - M.Tech CSE (AI), NSUT
%  Supervisor: Dr. Shobha Bhatt
%
%  Compile with XeLaTeX (Overleaf: Menu -> Compiler -> XeLaTeX).
%  Falls back to pdfLaTeX via the tipa guard below.
%  Requires nsut_logo.png in the same folder.
% =============================================================
\documentclass[12pt,a4paper]{article}

\usepackage{iftex}
\ifPDFTeX
  \usepackage[T1]{fontenc}
  \usepackage[utf8]{inputenc}
  \usepackage{tipa}
  \newcommand{\ipaschwa}{\textipa{@}}
  \newcommand{\ipalong}{\textipa{:}}
  \newcommand{\ipaA}[1]{/#1/}
\else
  \usepackage{fontspec}
  % \setmainfont{Charis SIL}   % uncomment if the default font lacks IPA
  \newcommand{\ipaschwa}{\symbol{"0259}}
  \newcommand{\ipalong}{\symbol{"02D0}}
  \newcommand{\ipaA}[1]{/#1/}
\fi

\usepackage[margin=1in]{geometry}
\usepackage{graphicx}
\usepackage{booktabs}
\usepackage{amsmath}
\usepackage{enumitem}
\usepackage{setspace}
\usepackage[hidelinks]{hyperref}
\usepackage[font=small,skip=4pt]{caption}
\usepackage{titlesec}

\onehalfspacing
\setlength{\parskip}{3pt}
\renewcommand{\arraystretch}{1.05}
\titlespacing*{\section}{0pt}{10pt}{4pt}
\titlespacing*{\subsection}{0pt}{7pt}{3pt}

\begin{document}

% --------------------------------------------------------------
% TITLE PAGE
% --------------------------------------------------------------
\begin{titlepage}
    \centering
    \vspace*{0.3cm}

    {\Large \textbf{Empirical Benchmarking of Neural Text-to-Speech\\
    Architectures for Indic Languages:\\
    Does Explicit Devanagari Grapheme-to-Phoneme\\
    Conversion Still Matter?}} \\[1.1cm]

    This progress report submitted in partial fulfillment of the requirements for the Degree of \\[0.5cm]

    {\large \textbf{MASTER OF TECHNOLOGY}} \\
    in \\
    Computer Science and Engineering - Artificial Intelligence \\[1.1cm]

    By \\[0.2cm]
    \textbf{Kaustubh Pandey} \hfill \textbf{2025PA17326} \\[1.1cm]

    Under the Supervision \\
    of \\[0.2cm]
    \textbf{Dr. Shobha Bhatt} \\
    \textbf{(Assistant Professor)} \\[1.1cm]

    \includegraphics[width=3.8cm]{nsut_logo.png} \\[1.1cm]

    Department of Computer Science and Engineering \\
    Netaji Subhas University of Technology (NSUT) \\
    New Delhi, India-110078 \\
    SEPTEMBER 2026
\end{titlepage}

\tableofcontents
\thispagestyle{empty}
\newpage
\setcounter{page}{1}

% ==============================================================
\section{Introduction}

\subsection{Background}

Two Hindi words make the problem concrete. \textit{Namak} (salt) is pronounced
\ipaA{n\ipaschwa m\ipaschwa k}; \textit{namkeen} (savoury) is
\ipaA{n\ipaschwa mki\ipalong n}. Both are written in Devanagari with no vowel
sign after the second consonant, yet one keeps a vowel there and the other
drops it. The spelling does not tell you which.

This is schwa deletion. Devanagari is an abugida, so every consonant carries an
unwritten inherent vowel, the schwa \ipaschwa, and Hindi deletes many of them.
The standard treatment is still the iterative right-to-left rule of Narasimhan,
Sproat and Kiraz \cite{narasimhan2004schwa}, which deletes a schwa lying
between two consonants that are themselves flanked by vowels. It works well but
not perfectly, with errors clustering at morpheme boundaries and in loanwords.
Arora et~al. \cite{arora2020schwa} later recast the task as supervised sequence
labelling and improved on it.

Synthesis itself has moved on entirely. Three architectures cover most of
current practice. FastSpeech~2 \cite{ren2021fastspeech2} is a
non-autoregressive transformer and it needs explicit durations. VITS
\cite{kim2021vits} maps text to waveform through a variational autoencoder
trained adversarially, and it learns its own alignment. Matcha-TTS
\cite{mehta2024matcha} uses conditional flow matching to reach
diffusion-level quality in far fewer steps. Where a spectrogram is produced,
HiFi-GAN \cite{kong2020hifigan} turns it into audio.

Indic resources have improved too. They run from the IIT Madras consortium
recordings \cite{baby2016resources} through AI4Bharat's baselines
\cite{kumar2023nextbillion} and IndicVoices-R \cite{sankar2024indicvoicesr} to
MMS \cite{pratap2024mms}. The MMS pretrained VITS checkpoints cover both my
languages, so fine-tuning is now the realistic starting point. One further fact
shapes this study. Marathi uses the same script as Hindi and shares most of its
phoneme inventory \cite{ohala1999hindi, dhongde2009marathi}, but it does not
apply Hindi's medial schwa deletion rule.

\subsection{Challenges}

Devanagari underspecifies the vowel, so a phonemic front end must guess, and
sometimes guesses wrong. Beyond that, published comparisons rarely hold
anything constant: two systems typically differ in architecture, data, sampling
rate, vocoder and front end at once, so no quality gap can be attributed to any
one of them.

The metrics are easy to get wrong too. MCD that keeps the zeroth coefficient
is partly a loudness measure \cite{kubichek1993mcd}. Pitch error on a linear
scale is dominated by the higher-pitched speaker. And word error rate means
nothing on its own. You need the recogniser's own error rate on natural audio,
which is the floor, and it is rarely reported.

\subsection{Research Gap}

Strong single-system Indic results exist, as do architectural comparisons for
English. What does not exist is a controlled measurement of how much the front
end contributes, separated from the architecture, on a Devanagari language.

Three things are missing specifically. No study varies only the input
representation while holding architecture, data, budget, vocoder and evaluation
fixed. No study uses a same-script control, so a difference found on Hindi
alone cannot distinguish "Devanagari is hard" from "schwa deletion is hard".
And no study reports how the effect changes with data quantity, though a team
with ten hours and a team with ten minutes face different questions.

\subsection{Motivation}

The practical question is whether to write a G2P front end at all. It is real
work, and building mine took most of a month. If
current architectures learn schwa deletion from data, that month was wasted and
nobody else should spend it. If they do not, the field should stop defaulting
to character-level input for Indic languages, which is presently common. Either
answer is useful and neither has been published.

\subsection{Problem Statement}

Holding training budget, evaluation protocol and speaker fixed, does phonemic
input from an explicit Devanagari G2P front end improve neural TTS quality in
Hindi relative to raw grapheme input? Does the answer depend on the acoustic
architecture? Does it survive as training data falls from nine hours to ten
minutes? And using Marathi as a same-script control where medial schwa deletion
does not apply, is any effect attributable to schwa deletion rather than to the
script in general?

\subsection{Research Objectives}

\begin{enumerate}[leftmargin=*, label=O\arabic*., itemsep=1pt, topsep=2pt]
\item Build a speaker-matched, loudness-normalised Hindi and Marathi corpus
with frozen, checksummed splits that regenerate identically.
\item Implement a Devanagari G2P front end with schwa deletion as a separately
switchable stage, and measure its accuracy against elicited speaker judgement
before trusting it downstream.
\item Train FastSpeech~2, VITS and Matcha-TTS under a provably identical
budget, on phonemic and on graphemic input.
\item Measure the phoneme-versus-grapheme effect across a nested data ladder
from nine hours to ten minutes.
\item Repeat the ablation in Marathi to separate phonology from script.
\item Evaluate with verified objective metrics, predicted MOS, a human
listening test, and real-time factor on consumer hardware.
\item Release the testbed, the checkpoints and a speaker-validated
contested-schwa stress-test set.
\end{enumerate}

% ==============================================================
\section{Literature Survey}

\textbf{Architectures.} FastSpeech~2 \cite{ren2021fastspeech2} (27M) predicts
duration, pitch and energy explicitly then generates the spectrogram in
parallel, which makes it fast but dependent on ground-truth durations from
somewhere, traditionally an external aligner. VITS \cite{kim2021vits} (83M)
collapses the pipeline into one model and learns alignment internally by
monotonic alignment search, inherited from Glow-TTS \cite{kim2020glowtts}.
Matcha-TTS \cite{mehta2024matcha} (18M) uses optimal-transport conditional flow
matching to match diffusion quality in fewer steps, which makes it interesting
at the low-data end of my ladder. StyleTTS~2 \cite{li2023styletts2} reports
stronger naturalness than any of them; I cut it for compute.

\textbf{Vocoding.} HiFi-GAN \cite{kong2020hifigan} remains the default. At
around 14M parameters for the V1 generator it is cheap beside the acoustic
models, which is why my design holds it fixed across the two systems needing
one.

\textbf{Alignment.} The Montreal Forced Aligner \cite{mcauliffe2017mfa} has
been the conventional duration source for FastSpeech-family models and is what
my original plan assumed. Badlani et~al. \cite{badlani2022alignment} then
showed that an alignment module trained jointly with the TTS model, without
external supervision, matches or beats forced-aligner durations. That result
changed my design; Section 4.5 explains why.

\textbf{Indic resources.} The consortium recordings \cite{baby2016resources}
sit behind most Indian-language TTS work, including both corpora here. Kumar
et~al. \cite{kumar2023nextbillion} showed usable quality from modest
per-language data. IndicVoices-R \cite{sankar2024indicvoicesr} scaled that up
for synthesis. MMS \cite{pratap2024mms} supplies the checkpoint my VITS
condition initialises from. Between them these papers establish that Indic TTS works and
how well it works. The contribution of the front end is not separated out in
any of them, and that is what I am measuring.

\textbf{Evaluation.} Mel-cepstral distortion \cite{kubichek1993mcd} is the
standard spectral measure; pitch and aperiodicity usually come from WORLD
\cite{morise2016world}; intelligibility is increasingly assessed by
transcribing synthesised audio with Whisper \cite{radford2023whisper}. Listening
tests are slow, so automatic MOS predictors such as UTMOS
\cite{saeki2022utmos} and NISQA \cite{mittag2021nisqa} have become common.
Both were trained largely on English and Japanese. I will measure how well they
correlate with my own human ratings and report that as a result. Loudness
normalisation follows ITU-R BS.1770 \cite{itu2015bs1770}.

% ==============================================================
\section{Experimental Setup and Work Done Till Now}

\textbf{Corpora and standardisation.} I use the Hindi and Marathi corpora from
the SPRING Lab IndicTTS collection. Neither dataset card states hours, speaker
count or sampling rate, so I measured all three. An early estimate from a
50-utterance sample proved 19\% low for Hindi. Every utterance then goes through the same chain. It is downmixed to mono,
silence-trimmed at 40~dB below peak with 50~ms of padding kept, and normalised
to $-23$~LUFS \cite{itu2015bs1770}. From there it is resampled twice, to
22.05~kHz for the master and separately to 16~kHz for the VITS arm, and
peak-guarded at 0.99. Two details in that chain are deliberate. Trimming happens at the 48~kHz
source, so both outputs share their boundaries to the sample. And the 16~kHz
copy comes from that same gained source, not from the master. Cascading two
resamples would put a measurably different signal in front of VITS than the
other architectures see, and that difference would show up in the results
looking like an architectural effect.

\textbf{G2P front end.} Four separable stages, arranged so the phonological
rule can be switched without touching anything else. Normalisation handles NFC,
nukta composition, corpus spelling repair, digit expansion and punctuation
folding. Segmentation makes every inherent schwa explicit and tags it as
unwritten. Schwa deletion applies the rule from \cite{narasimhan2004schwa} plus
obligatory word-final deletion plus a lexicon override; Marathi runs identical
code with medial deletion disabled, and that single flag is the whole control
condition. Finally everything maps into one IPA inventory shared by both
languages. A parallel grapheme path reuses the identical normalisation chain,
so the ablation measures phonology rather than preprocessing.

\textbf{Gold pronunciations without circularity.} I wrote the schwa rule, so
gold forms I also wrote would measure only my own consistency. They come from
native speakers instead, through a sheet that never shows the rule's
prediction. Each word appears with its candidate vowels spelled out in numbered
slots, as in \texttt{n(a1)m(a2)kiin(a3)}. Speakers mark each one yes, no or
variable. They do not write IPA, because transcription by non-phoneticians adds
noise that has nothing to do with the question. Their answers pass through the same converter
to produce IPA. Items where speakers disagree are marked variable, excluded
from the accuracy denominator and reported as a variation rate.

\textbf{Metrics.} MCD uses DTW alignment over 24 coefficients, excluding
$c_0$. The alignment matters: a model with a duration predictor differs
systematically in timing from an end-to-end one, and frame-aligned comparison
would charge that timing difference as spectral error. Prosody is reported as
two separate numbers, log-$F_0$ RMSE over commonly voiced frames and
voiced/unvoiced error rate. A system with correct voicing and wrong pitch has a
different defect from one with the reverse.

\textbf{Design and reproducibility.} There are seventeen configurations
(Table~\ref{tab:runs}). Every acoustic run shares 100{,}000 steps, 12{,}000 mel
frames per batch, a learning rate of $2\times10^{-4}$ with inverse-square-root
decay after 4{,}000 warmup steps, fp16, and gradient clipping at 1.0. That
sentence is the study's central claim, so an assertion enforces it. I am not
relying on remembering it. One YAML file defines one run, and it carries a hash
over every field that affects the trained model. A second assertion refuses two
runs that share a hash. The splits are locked with a SHA-256 per file.
Preparation, G2P and evaluation all run locally on an Apple M5. Training is
planned on free-tier GPU resources, capped at 30 GPU-hours a week in 12-hour
sessions, which is where the checkpoint-resume requirement comes from.

% ==============================================================
\section{Results and Discussion}

Nothing has been trained yet, so what follows are properties of the corpus, the
front end and the measuring apparatus. Each had to be established before any
model was trained. Several turned out to be findings in their own right.


\subsection{Corpus Profile}

Neither dataset labels its speakers; there is a binary field with no
explanation. Hindi and Marathi have to be matched on speaker
characteristics, so I recovered the labels from median $F_0$ over 25 sampled
utterances per speaker. It separates them cleanly, and identically, in both
languages (Table~\ref{tab:corpus}). Every gender label here rests on that
measurement. The dataset itself says nothing about it.

\begin{table}[htbp]
\centering\small
\caption{Measured corpus properties. Gender is inferred from median $F_0$.}
\label{tab:corpus}
\begin{tabular}{lrrrr}
\toprule
& \multicolumn{2}{c}{Hindi} & \multicolumn{2}{c}{Marathi} \\
\cmidrule(lr){2-3}\cmidrule(lr){4-5}
& Speaker 0 & Speaker 1 & Speaker 0 & Speaker 1 \\
\midrule
Utterances        & 5{,}983 & 5{,}842 & 5{,}362 & 5{,}577 \\
Hours (raw)       & 13.10 & 13.13 & 11.49 & 10.88 \\
Median $F_0$ (Hz) & 198.9 & 108.8 & 237.2 & 125.7 \\
Inferred gender   & female & male & female & male \\
\midrule
Corpus total & \multicolumn{2}{c}{11{,}825 utts, 26.23 h, 48 kHz} &
               \multicolumn{2}{c}{10{,}939 utts, 22.37 h, 48 kHz} \\
\bottomrule
\end{tabular}
\end{table}

I selected the male speaker in both. The hours are near identical either way, so pitch
decided it. The two male voices sit 17~Hz apart across the two languages. The
female pair sit 38~Hz apart. Picking the men leaves less acoustic difference
between the arms that has nothing to do with schwa deletion.

I checked transcript quality with characters per second of audio. A transcript
that belongs to different audio shows up as an outlier on that measure. Hindi's
median is 12.52 (p1 9.16, p99 15.77) and Marathi's is 9.10 (6.64, 11.70).
Nothing in either falls outside 0.4 to 2.0 times the median, so there is
nothing to clean. The longest Hindi utterance runs 127.7~s with 1{,}600
characters. I assumed it was corrupt. That works out at 12.5 characters per
second, which is exactly the corpus median.

\subsection{Standardisation and Splits}

Trimming removed 5.5\% of Hindi and 8.4\% of Marathi. Marathi carried
noticeably more dead air. That had a consequence I had not planned. After
trimming, and after the frozen test and development sets are taken out, only
9.53~h and 9.16~h of training audio survive. My ladder was going to start at
ten hours. It now starts at nine, because matching the rungs across the two
languages matters more than a round number. An unmatched top rung would drop a
data-quantity difference into the control comparison, and that is the one place
it must not be.

\begin{table}[htbp]
\centering\small
\caption{Standardised corpora (male speaker, 1--15\,s, $-23$\,LUFS) and the
nested data ladder in utterances per rung.}
\label{tab:prepared}
\begin{tabular}{lrr@{\hskip 20pt}lrr}
\toprule
& Hindi & Marathi & Rung & Hindi & Marathi \\
\midrule
Utterances retained & 5{,}485 & 5{,}559 & 9 h     & 4{,}793 & 5{,}071 \\
Hours               & 10.289 & 9.884   & 5 h     & 2{,}668 & 2{,}810 \\
Peak-limited        & 1 & 0             & 1 h     & 541 & 562 \\
Test / dev          & 300 / 100 & 300 / 100 & 30 min & 265 & 279 \\
Train (hours)       & 5{,}085 (9.525) & 5{,}159 (9.160) & 10 min & 86 & 90 \\
\bottomrule
\end{tabular}
\end{table}

Spot-checking confirmed exactly $-23.00$~LUFS at both rates. The two copies
agree on duration to under a millisecond, the peaks stay under the ceiling, and
the edge padding is there. The peak guard fired once. So it is not dead code.

Split membership comes from a salted SHA-256 of the utterance identifier. A
seeded shuffle, the usual choice, depends on the order the manifest happened to
be written in, so re-running preparation elsewhere can hand you a quietly
different test set. Hashing the identifier makes membership a property of the
utterance, and rebuilding reproduces the test split byte for byte. The ladder
takes prefixes of that one ordering. Sampling each rung independently would let
sampling variance look like a data-quantity effect, which is the effect the
ladder exists to measure. I verified that every rung is a strict
superset of the one below, that the three splits are pairwise disjoint, and
that no rung reaches outside training.

\subsection{G2P Coverage}

I developed the front end against 49 hand-picked words and it handled all of
them. Then I ran it over all 11{,}044 corpus transcripts. 4.1\% of Hindi phone
tokens and 2.7\% of Marathi had no mapping at all. Both are now zero (Table~\ref{tab:coverage}). Getting there meant fixing seven
things a curated test set was never going to find. Nine Hindi tokens carried a
nukta on a consonant that has no nukta form, and the stray mark also broke the
parse of the next character. 199 Marathi tokens were the letter RRA. There were
keyboard slips producing an invalid candra-O sequence, and a colon typed where
a visarga belonged. The OM ligature is a syllable written as a single letter.
Southern short vowels appear that neither language distinguishes. And some
vowel signs had no consonant to attach to.

\begin{table}[htbp]
\centering\small
\caption{G2P coverage over both corpora after correction.}
\label{tab:coverage}
\begin{tabular}{lrr}
\toprule
& Hindi & Marathi \\
\midrule
Word tokens & 99{,}958 & 55{,}492 \\
Phone tokens & 494{,}863 & 375{,}684 \\
Unknown symbols & 0 & 0 \\
Inventory used / declared & 72 / 75 & 67 / 75 \\
\bottomrule
\end{tabular}
\end{table}

Marathi's inventory is a strict subset of Hindi's. It lacks only five marginal
or Perso-Arabic phones, so the two arms share one symbol table with no
Marathi-only embedding rows. One weakness is worth recording here. Three
Marathi phones occur fewer than three times in the whole corpus, so those
embedding rows will get essentially no gradient.

\subsection{Metric Verification}

I checked both metrics against answers I could derive on paper.

MCD proved amplitude invariant. Scaling the input by 0.25 changes it by under
$10^{-6}$~dB. That is the property that confirms the implementation. Scaling
multiplies the power spectrum by a constant, which is a constant decibel offset
across every mel band, and that lands entirely in the zeroth DCT coefficient.
Dropping $c_0$ therefore makes the measure blind to gain. If that test ever
failed, every MCD figure here would be partly a loudness measurement, and
nothing in the results would show it.

Log-$F_0$ RMSE reports $\ln(1.1)$ for a 10\% pitch error at 100~Hz and again
at 300~Hz, exact to $10^{-12}$. A linear RMSE would call the same perceptual
error 10~Hz on one voice and 30~Hz on another. Two further guards are tested.
Unequal-length $F_0$ tracks raise rather than truncate, because truncation
would favour systems that predict shorter durations. And a comparison with no
commonly voiced frames returns NaN, not zero. The suite stands at 78 tests.

\subsection{Three Defects Found}

Each would have corrupted results without any visible symptom, and two were
mine.

\textbf{Punctuation was rewriting words.} I kept punctuation as a prosodic
token, since it is the only phrasing cue the text carries. Leaving it attached
through the schwa stage breaks the rule twice over. A trailing comma means the
word-final schwa is no longer final, so it survives. Its survival then supplies
the right-hand vowel that lets the schwa to its left delete. \textit{Namak},
\ipaA{n\ipaschwa m\ipaschwa k}, came out \ipaA{n\ipaschwa mk\ipaschwa}, a
different word. 16.5\% of Hindi word tokens and 14.9\% of Marathi carry
punctuation. The damage would have fallen on the phonemic arm alone, because
the grapheme path never touches the schwa stage. My headline comparison would
have been partly measuring this bug. Punctuation is now stripped before the
schwa stage and restored after. There is an invariance test across nine words
and four marks in both languages.

\textbf{A run was in the plan twice.} The design had 19 runs. The ladder's
nine-hour rung was a second copy of the main run under a different identifier.
That is ten GPU-hours for a number I already had. It stayed invisible because
the run identifier sat inside the configuration hash, which by construction
makes two identical runs look distinct. The identifier is now excluded, and an
assertion refuses any two runs that share a settings hash.

\textbf{The elicitation sheet gave away its own answers.} Its worked examples
were themselves items from the 49-word test set. So every participant was
handed the answers to three of the words they were about to be tested on. I
noticed only because someone reading the sheet asked about one. Examples now come from outside the
set, enforced by a check at generation time.

\subsection{Two Design Decisions}

\textbf{Forced alignment dropped.} My plan took FastSpeech~2's durations from
MFA \cite{mcauliffe2017mfa}. I have replaced that with internal unsupervised
alignment \cite{badlani2022alignment}, for three reasons. The first is that it
was an uncontrolled asymmetry. VITS and Matcha both learn alignment internally
\cite{kim2020glowtts}, so externally supervised durations would give one
architecture information the other two never receive, inside a comparison whose
subject is the architecture. The second is that it was asymmetric across
languages as well. MFA publishes a pretrained Hindi model and no Marathi one,
which puts that difference inside the control comparison. The third is that MFA
would not run in my environment. The package installs and then fails on Kaldi
bindings, which have no wheel for my architecture.

One limitation follows from the change. Internal alignment is weaker when data
is very scarce, so the lowest ladder rungs may suffer more than the top ones. I
will measure that as a property of the ladder and report it.

\textbf{Sampling rate.} The master is 22.05~kHz, as the synopsis specified.
But VITS initialises from an MMS checkpoint \cite{pratap2024mms} that is
natively 16~kHz, so its output carries nothing above 8~kHz. Full-band MCD would
then charge it for missing bandwidth rather than for worse modelling. MCD is reported in two columns, full band and band-limited to
8~kHz, and the band-limited column is the headline. Listening stimuli are
low-passed to 8~kHz for every system, so raters hear equal bandwidth. ASR-WER
needs no mitigation at all, since Whisper resamples to 16~kHz anyway
\cite{radford2023whisper}.

\subsection{Compute Requirements}

\begin{table}[htbp]
\centering\small
\caption{The experimental matrix. All acoustic runs share one budget, enforced
by assertion.}
\label{tab:runs}
\begin{tabular}{llllrr}
\toprule
Run & Architecture & Lang. & Input & Rate (kHz) & GPU-h \\
\midrule
r01 / r02 / r03 & FS2 / VITS / Matcha & hi & phoneme & 22.05 / 16 / 22.05 & 36 \\
r04 / r05 & FS2 / VITS & hi & grapheme & 22.05 / 16 & 16 \\
r06 & HiFi-GAN & hi & n/a & 22.05 & 10 \\
r07--r10 & FastSpeech 2 & hi & phoneme, ladder & 22.05 & 20 \\
r11--r14 & VITS & hi & phoneme, ladder & 16 & 20 \\
r15 / r16 & FS2 / VITS & mr & phoneme & 22.05 / 16 & 24 \\
r17 & HiFi-GAN & mr & n/a & 22.05 & 10 \\
\midrule
\multicolumn{5}{r}{\textbf{Total (17 runs)}} & \textbf{136} \\
\bottomrule
\end{tabular}
\end{table}

Table~\ref{tab:runs} totals 136 GPU-hours at one seed per configuration. On
free-tier resources capped at 30 hours a week, that is 4.53 weeks of training
alone. One seed is not really enough. My rule is that no difference smaller
than the measured seed variance gets reported as a finding, and with a single
seed everywhere that floor rests on replicating one cell and generalising from
it. Three seeds throughout would take the requirement to 368 GPU-hours.

Institutional GPU access would mostly buy statistical strength. On a DGX A100
the single-seed campaign finishes in roughly 1.6 days, and three-seed
replication in about 4.4. The calendar saving is around four weeks out of
thirteen, since evaluation and the listening study wait on people. The larger
gain is in what I can claim. Every comparison moves from ``these numbers
differ'' to ``these differ by more than run-to-run noise''.

% ==============================================================
\section{Conclusion}

Preparation is done. Two Devanagari corpora totalling 48.6 hours have been
profiled, standardised, and split into frozen checksummed partitions that
regenerate identically. Matching them on speaker characteristics meant
recovering those characteristics acoustically first, because the datasets do
not record them. The G2P front end is built, with schwa deletion isolated as a
switchable stage, and it covers both corpora completely. The instrument for
validating it against native-speaker judgement is waiting on participants. Both
objective metrics are checked against analytic answers. All 17 runs are
specified, and the shared budget is enforced by assertion.

Three defects turned up along the way and none of them announced itself. The
worst affected 16.5\% of Hindi word tokens and would have degraded only the
phonemic arm, which is the half of the comparison this study rests on. All
three were found the same way: by running the code over the real corpus, and by
watching someone try to use the elicitation sheet. The 78 unit tests caught
none of them. I will pilot the listening study on two or three people before
recruiting twenty.

Remaining immediately: collecting elicitation sheets, confirming the GPU
execution path end to end, and writing the training wrappers.

% ==============================================================
\section*{References}
\addcontentsline{toc}{section}{References}
\begingroup
\footnotesize
\setstretch{1.0}
\begin{thebibliography}{99}
\setlength{\itemsep}{0pt}\setlength{\parsep}{0pt}

\bibitem{ren2021fastspeech2}
Y. Ren et~al., ``FastSpeech 2: Fast and high-quality end-to-end text to
speech,'' \textit{ICLR}, 2021. arXiv:2006.04558.

\bibitem{kim2021vits}
J. Kim, J. Kong, and J. Son, ``Conditional variational autoencoder with
adversarial learning for end-to-end text-to-speech,'' \textit{ICML}, 2021.
arXiv:2106.06103.

\bibitem{mehta2024matcha}
S. Mehta et~al., ``Matcha-TTS: A fast
TTS architecture with conditional flow matching,'' \textit{ICASSP}, 2024.
arXiv:2309.03199.

\bibitem{kong2020hifigan}
J. Kong, J. Kim, and J. Bae, ``HiFi-GAN: Generative adversarial networks for
efficient and high fidelity speech synthesis,'' \textit{NeurIPS}, 2020.
arXiv:2010.05646.

\bibitem{kim2020glowtts}
J. Kim et~al., ``Glow-TTS: A generative flow for
text-to-speech via monotonic alignment search,'' \textit{NeurIPS}, 2020.
arXiv:2005.11129.

\bibitem{li2023styletts2}
Y. A. Li et~al., ``StyleTTS 2:
Towards human-level text-to-speech through style diffusion and adversarial
training with large speech language models,'' \textit{NeurIPS}, 2023.
arXiv:2306.07691.

\bibitem{badlani2022alignment}
R. Badlani et~al.,
``One TTS alignment to rule them all,'' \textit{ICASSP}, 2022.
arXiv:2108.10447.

\bibitem{mcauliffe2017mfa}
M. McAuliffe et~al., ``Montreal
Forced Aligner: Trainable text-speech alignment using Kaldi,''
\textit{Interspeech}, 2017, pp. 498--502.

\bibitem{pratap2024mms}
V. Pratap et~al., ``Scaling speech technology to 1,000+ languages,''
\textit{Journal of Machine Learning Research}, vol. 25, 2024. arXiv:2305.13516.

\bibitem{kumar2023nextbillion}
G. K. Kumar et~al., ``Towards
building text-to-speech systems for the next billion users,'' \textit{ICASSP},
2023. arXiv:2211.09536.

\bibitem{sankar2024indicvoicesr}
A. Sankar et~al., ``IndicVoices-R: Unlocking a massive multilingual
multi-speaker speech corpus for scaling Indian TTS,'' \textit{NeurIPS Datasets
and Benchmarks}, 2024. arXiv:2409.05356.

\bibitem{baby2016resources}
A. Baby, A. L. Thomas, N. L. Nishanthi et~al., ``Resources for
Indian languages,'' \textit{Proc. Text, Speech and Dialogue (CBBLR Workshop)},
2016.

\bibitem{narasimhan2004schwa}
B. Narasimhan, R. Sproat, and G. Kiraz, ``Schwa-deletion in Hindi text-to-speech
synthesis,'' \textit{International Journal of Speech Technology}, vol. 7,
no. 4, pp. 319--333, 2004.

\bibitem{arora2020schwa}
A. Arora, S. Gessler, and N. Hundt, ``Supervised grapheme-to-phoneme conversion
of orthographic schwas in Hindi and Punjabi,'' \textit{ACL}, 2020.
arXiv:2004.10353.

\bibitem{ohala1999hindi}
M. Ohala, ``Hindi,'' in \textit{Handbook of the International Phonetic
Association}, Cambridge University Press, 1999, pp. 100--103.

\bibitem{dhongde2009marathi}
R. V. Dhongde and K. Wali, \textit{Marathi}, John Benjamins, 2009.

\bibitem{kubichek1993mcd}
R. Kubichek, ``Mel-cepstral distance measure for objective speech quality
assessment,'' \textit{IEEE Pacific Rim Conf. on Communications, Computers and
Signal Processing}, 1993, pp. 125--128.

\bibitem{morise2016world}
M. Morise, F. Yokomori, and K. Ozawa, ``WORLD: A vocoder-based high-quality
speech synthesis system for real-time applications,'' \textit{IEICE Trans. on
Information and Systems}, vol. E99-D, no. 7, pp. 1877--1884, 2016.

\bibitem{radford2023whisper}
A. Radford et~al.,
``Robust speech recognition via large-scale weak supervision,'' \textit{ICML},
2023. arXiv:2212.04356.

\bibitem{saeki2022utmos}
T. Saeki et~al., ``UTMOS: UTokyo-SaruLab system for VoiceMOS Challenge 2022,''
\textit{Interspeech}, 2022. arXiv:2204.02152.

\bibitem{mittag2021nisqa}
G. Mittag et~al., ``NISQA: A deep
CNN-self-attention model for multidimensional speech quality prediction with
crowdsourced datasets,'' \textit{Interspeech}, 2021. arXiv:2104.09494.

\bibitem{itu2015bs1770}
International Telecommunication Union, ``Algorithms to measure audio programme
loudness and true-peak audio level,'' Recommendation ITU-R BS.1770-4, 2015.

\end{thebibliography}
\endgroup

\end{document}
